\section{研究の要約}

本研究では、エッジ環境（Webブラウザ）における待機列のグループ検出精度向上およびリアルタイム処理の実現に向け、段階的な13のアプローチを通して技術的検証を行ってきた。

各アプローチにより得られた主な知見および達成事項は以下の通りである。

\begin{enumerate}
    \item \textbf{前処理・検出領域の最適化（アプローチ01〜02）}：
    二次元距離のみに依存するベースライン手法（アプローチ01）の誤検出に対し、関心領域（ROI）の空間制限（アプローチ02）を導入することで背景ノイズや周囲通行人の干渉を防止した。
    \item \textbf{検出対象の試行と姿勢情報の導入（アプローチ03〜04）}：
    顔検出（アプローチ03）および骨格・ポーズ推定（アプローチ04）の導入により、密集部位での遮蔽問題や後ろ姿・横顔に対する検出漏れを段階的に解消した。
    \item \textbf{ハイブリッド構成とパフォーマンス最適化（アプローチ05〜08）}：
    一次検出とポーズ推定を組み合わせたパイプライン（アプローチ05）に対し、Worker内サイズ保証（アプローチ06）、IoUによる重複排除（アプローチ07）、縦方向の断片化結合およびWebGPU/16ビット精度によるメモリ最適化（アプローチ08）を実施し、エッジでの安定動作を実現した。
    \item \textbf{幾何構造と検出エンジンの刷新（アプローチ09〜10）}：
    身体向きベクトルと列線計算（アプローチ09）により蛇行列に対応し、さらに ONNX Runtime Web 上で動作する YOLOv12 Nano への全面刷新およびWorker一本化（アプローチ10）により一次検出の精度とスピードを大幅に向上させた。
    \item \textbf{対向関係検出と動的トラッキング（アプローチ11〜13）}：
    身体向きベクトルの延長線交差判定（アプローチ11）により前後の列に跨がる対話グループの検出能力を獲得し、時空間情報とカルマンフィルタを組み合わせたトラッキング（アプローチ12〜13）により、軽量な計算資源のままで一時的な通行人の排除を可能とした。
\end{enumerate}

以上の取り組みにより、高負荷なサーバー環境に依存することなく、ブラウザエッジ環境上で動的かつ高精度に待機列とグループ構成を推定するシステム設計が確立された。
